\documentclass[12pt]{article}
\usepackage[utf8]{inputenc}
\usepackage[brazil]{babel}
\usepackage{lmodern}
\usepackage{enumitem}
\usepackage{geometry}
\geometry{margin=2.5cm}

\usepackage{amssymb,amsmath,amsfonts,eurosym,geometry,ulem,graphicx,caption, xcolor,setspace,sectsty,comment,footmisc,caption,natbib,pdflscape,array,hyperref, fancyhdr, csquotes, tikz, enumitem, dsfont, amsthm, makecell, multicol,multirow, booktabs,xcolor,subcaption}

\title{Guia -- Matching -- Outline}
\author{}
\date{}

\begin{document}

\maketitle
\doublespacing

\section{Introdução: Por que Precisamos de Comparações Justas}
A pergunta central em avaliação de políticas é simples: \textbf{``Esse programa funciona?''}. 
No entanto, responder a essa questão exige comparações justas entre beneficiários e não beneficiários. 
Se os participantes de um curso de qualificação, por exemplo, já eram mais motivados ou tinham melhor escolaridade antes mesmo do curso, comparar médias de resultados diretamente pode levar a conclusões enviesadas. 
Esse é o problema clássico do \textit{viés de seleção}.

\section{A Ideia Central do Matching}
O matching busca resolver esse problema aproximando-se de uma situação contrafactual: encontrar um grupo de comparação que seja o mais parecido possível com os participantes do programa. 
Diferente de simplesmente comparar médias, o matching tenta construir um grupo de controle que reflita como seriam os participantes caso não tivessem recebido o tratamento.

\section{Matching Exato}
No matching exato, comparamos indivíduos que compartilham exatamente as mesmas características observáveis relevantes (por exemplo, idade, escolaridade e gênero). 
Isso pode ser ilustrado com uma tabela simples, em que cada tratado é pareado com um controle idêntico.  

\textbf{Prós:} é altamente intuitivo e transparente.  
\textbf{Contras:} com muitas variáveis, torna-se praticamente inviável devido à chamada \textit{maldição da dimensionalidade}.

\section{Matching Aproximado}
Na prática, o matching exato geralmente falha porque raramente encontramos pares perfeitamente idênticos. 
Por isso, métodos aproximados são amplamente usados:

\begin{itemize}[topsep=2pt,itemsep=2pt]
    \item \textbf{Propensity Score Matching:} condensa várias variáveis observadas em uma única medida de probabilidade de tratamento.
    \item \textbf{Nearest-Neighbor Matching:} seleciona o indivíduo mais próximo em termos de características.
    \item \textbf{Kernel Matching / Ponderação:} constrói médias ponderadas de vários controles.
    \item \textbf{Coarsened Exact Matching:} agrupa variáveis em categorias mais amplas para permitir pareamentos.
\end{itemize}

Essas abordagens podem ser ilustradas com gráficos que mostrem como controles são ``trazidos para perto'' dos tratados.

\section{Matching e Outros Métodos Econométricos}
O matching não é uma técnica isolada: muitas ferramentas modernas de inferência causal incorporam sua lógica.

\begin{itemize}[topsep=2pt,itemsep=2pt]
    \item \textbf{Matching + Diferenças-em-Diferenças (DiD):} ajuda a garantir que grupos tratados e de controle sejam mais comparáveis antes da análise de trajetórias temporais.
    \item \textbf{Controle Sintético (SC):} cria uma ``unidade sintética'' a partir de combinações de várias unidades de controle — essencialmente, um matching em nível agregado.
    \item \textbf{DiD Sintético:} une elementos de DiD e SC, combinando matching e análise de tendências.
\end{itemize}

Conclusão parcial: as ideias de matching estão no centro de várias ferramentas modernas.

\section{Matching em RCTs}
Mesmo em experimentos randomizados, matching pode ter papel importante.  
Dois exemplos comuns:

\begin{itemize}[topsep=2pt,itemsep=2pt]
    \item \textbf{Pareamento (matched-pairs):} formar pares de indivíduos ou unidades semelhantes antes da randomização.
    \item \textbf{Estratificação:} dividir amostras em subgrupos equilibrados em características-chave.
\end{itemize}

Isso aumenta a precisão estatística e a credibilidade política dos resultados.

\section{O que Formuladores de Políticas Devem Ter em Mente}
O matching tem \textbf{pontos fortes}: é intuitivo, baseado em dados e transparente.  
Mas também apresenta \textbf{limitações}: não elimina viés de fatores não observados.  

Perguntas-chave:
\begin{itemize}[topsep=2pt,itemsep=2pt]
    \item Como os pares foram escolhidos?
    \item Tratados e controles eram realmente comparáveis?
    \item Podem existir diferenças ocultas que distorçam o resultado?
\end{itemize}

\section{Aplicações em Políticas Públicas}
Diversos estudos já usaram matching para avaliar programas em áreas como:
\begin{itemize}[topsep=2pt,itemsep=2pt]
    \item Educação (impacto de cursos técnicos ou bolsas de estudo).
    \item Saúde (programas de prevenção ou acompanhamento médico).
    \item Mercado de trabalho (políticas de qualificação ou subsídios).
\end{itemize}

Esses exemplos mostram como o matching contribui para avaliações mais convincentes.

\section{Principais Lições para Formuladores de Políticas}
\begin{itemize}[topsep=2pt,itemsep=2pt]
    \item Matching leva a comparações mais justas.
    \item Pode ser usado sozinho ou combinado com DiD, SC e RCTs.
    \item Aumenta a credibilidade, mas não garante causalidade absoluta.
\end{itemize}

\section*{Apêndice (Opcional, para Leitores Técnicos)}
Para leitores com interesse técnico, pode-se detalhar:
\begin{itemize}[topsep=2pt,itemsep=2pt]
    \item Definições formais de \textit{propensity score}.
    \item Propriedades assintóticas dos estimadores.
    \item Discussões sobre viés e variância.
\end{itemize}

\subsection*{Leituras Recomendadas}
\begin{itemize}[topsep=2pt,itemsep=2pt]
    \item Cunningham (2021). \textit{Mixtape}, Capítulo 5.
    \item Stuart (2010). \textit{Matching methods for causal inference: A review and a look forward}.
    \item Rosenbaum \& Rubin (1983, 1985).
    \item Abadie \& Imbens (2006, 2011).
\end{itemize}

\end{document}
